\section{Introduction}
Quantized edge inference exposes a discrete reliability problem: changing one stored bit can change a parameter by a small step or reverse its sign with a large numerical jump. Protecting every parameter bit may be unnecessary for a particular application, but protecting only a fraction requires a defensible allocation rule. A practical software study must therefore connect a ranking to the task quality it preserves, not merely demonstrate that its scores correlate with a proxy for sensitivity.

We investigate the question: \emph{How reliably can static encoded-bit susceptibility rankings guide budget-constrained protection under simultaneous parameter faults?} Three distinctions are central. First, a numerical bit jump is not task damage. Second, an estimated derivative is not the finite effect of an encoded XOR. Third, damage measured with a single faulty bit is not generally the benefit of suppressing that bit when other bits may also be faulty. A ranking can be informative at one stage and disappointing at the next.

NetSuRF addresses these distinctions through an end-to-end software evaluation flow. It binds a frozen quantized model and dataset to an exact eligible bit population; constructs susceptibility scores and protection masks; applies matched persistent fault fields; and reports task metrics and computational cost. The tested protection operation perfectly suppresses proposed events on selected coordinates. Its budget is the number or fraction of eligible bits selected. This is a controlled logical comparison, not an implementation or cost model for a physical mitigation mechanism. Throughout the paper, a reliability curve means application quality under the specified logical fault distribution, not a device failure rate.

Our earlier FCCM work introduced a selective-protection design flow and a qualitative gradient/delta-based technique~\cite{guo2025inpixel}. The present study does not claim that flow as new. It isolates and evaluates the \emph{software ranking-to-allocation decision}: frozen INT8 codes, corrected baselines, authoritative faulted inference, common random fields, explicit task metrics, and independent evaluation of bounded allocation choices. The published FKeras method is evaluated alongside significance and first-order baselines, preserving its tier-first allocation rule.

The contributions are threefold. (1) An encoding-aware evaluation methodology separates numerical fault semantics, surrogate scores, ordering policies, and measured protection outcomes. (2) A controlled comparison across image classification, keyword spotting, and anomaly detection measures isolated-damage association, task quality versus BER and coverage, and score-construction cost, including unfavorable results. (3) A compact analytical and empirical assessment explains why isolated rankings lack a universal allocation guarantee while quantifying how little current bounded-pool evidence establishes about improvement over an exact isolated reference. The underlying probability identities are established mathematical tools, not newly invented foundations.

The findings support conditional engineering guidance rather than a universal winner. Per-example absolute gradients substantially outperform MSB-first on some ResNet-8 conditions; keyword-spotting gains are smaller, and ToyCar results depend strongly on budget and task metric. A ranking that predicts reconstruction MSE is not thereby validated for ROC-AUC protection. Likewise, a numerically validated Hessian score cannot be credited with effectiveness when the tested budget produces the same mask as MSB-first. These distinctions define both the useful results and the limits of this internal study.
